\chapter{Reflection}
\label{ch:reflection}

This chapter details our personal reflections on this engineering project, discussing our decision-making process, the technical hurdles we resolved, and the key lessons we learned as a team.

% ================================================================

\noindent\textbf{Skills and Competence Developed.} \\
Before starting this project, our knowledge of reinforcement learning was mostly academic, based on the lectures given by Prof. Hirchoua Badr at ENSAM Casablanca. Setting up and training our agents on the 2v2 soccer task turned those equations into actual engineering skills. 

We now understand how the components of PPO (like the clipping ratio, advantage estimations, and entropy coefficients) work in practice to prevent policy collapse. Integrating the different libraries (Shimmy, PettingZoo, SuperSuit, and Stable-Baselines3) was an exercise in systems engineering. We had to debug dependency conflicts and resolve pickling issues that occurred when PyTorch tried to serialize our custom wrapper across multiple CPU cores.

\bigskip
\noindent\textbf{Problem Solving and Technical Hurdles.} \\
One major issue we faced was a Windows terminal encoding mismatch. Our python scripts output UTF-8 soccer pitch diagrams and logs, but the Windows command line expected CP1252 encoding. This caused the scripts to crash with encoding errors. We fixed this by modifying our scripts to reconfigure the standard output to UTF-8 using:
\begin{lstlisting}[language=Python]
import sys
if sys.platform.startswith('win'):
    sys.stdout.reconfigure(encoding='utf-8')
\end{lstlisting}
This taught us that small environment differences can cause major bugs in a deployment pipeline. 

Another challenge was our hardware limit: our local laptops lacked a dedicated NVIDIA GPU, making local training extremely slow. Moving our training to Google Colab allowed us to leverage a T4 GPU, but we had to manage the 12-hour session limits. We solved this by writing a custom callback that saved model checkpoints directly to our Google Drive every 200,000 steps, ensuring we never lost more than an hour of training progress.

\bigskip
\noindent\textbf{Teamwork and Development Strategy.} \\
Working together as a team (Imad Eddine Oukrati and Mohamed Echchyoughi), we divided our tasks: one of us focused on modeling the reward shaping wrapper and running training loops on Colab, while the other built the spatial tracking and analytics scripts to extract the 2D hexbin maps. 

If we were to start over, we would implement integration tests much earlier. We spent a lot of time debugging serialization errors late in the project. Running a short 1,000-step test run at the start would have saved us days of troubleshooting. We would also set up TensorBoard from the very first run to monitor convergence in real-time.

\bigskip
\noindent\textbf{Impact on our Engineering Approach.} \\
This project showed us the difference between reading about an algorithm in a textbook and implementing it in a physical simulator. The issues we encountered (multiprocessing serialization, encoding errors, and resource constraints) are rarely covered in theory, yet they represent a large portion of a machine learning engineer's daily work. This experience has prepared us to handle complex, continuous multi-agent systems in our future studies and careers.

% ================================================================
\section{Summary}
\label{sec:reflection_summary}
% ================================================================

This chapter concludes our report by outlining our personal learning journey and reflecting on our team's engineering choices. The References section follows below, detailing all academic papers and software libraries cited throughout our project.